\documentclass[../../../include/open-logic-section]{subfiles}

\begin{document}

\olfileid{sth}{ordinals}{opps}
\olsection{الأعداد الترتيبية الخلفية والحدية}

يكثر استعمال الأعداد الترتيبية في الفصول الآتية؛ فلنمهد لها ببعض
المفاهيم اليسيرة.

\begin{defn}
\emph{خلف} العدد الترتيبي~$\OLId{greek-variable}{alpha}$ هو~$\ordsucc{\OLId{greek-variable}{alpha}} =\OLId{greek-variable}{alpha} \cup \{\OLId{greek-variable}{alpha}\}$.
ويسمى~$\OLId{greek-variable}{alpha}$ عددًا ترتيبيًا \emph{خلفيًا} إذا كان
$\ordsucc{\beta} = \OLId{greek-variable}{alpha}$ لعدد ترتيبي~$\beta$. أما~$\OLId{greek-variable}{alpha}$ فيسمى
عددًا ترتيبيًا \emph{حديًا} إذا وفقط إذا كان غير خال وغير خلفي.
\end{defn}
\noindent
ووجه صحة هذا التعريف لـ\emph{الخلف} يظهر من القضية التالية.

\begin{prop}
لكل عدد ترتيبي~$\OLId{greek-variable}{alpha}$ تصح الأمور الآتية:
\begin{enumerate}
\item $\OLId{greek-variable}{alpha} \in \ordsucc{\OLId{greek-variable}{alpha}}$؛
\item $\ordsucc{\OLId{greek-variable}{alpha}}$ عدد ترتيبي؛
\item ليس من عدد ترتيبي~$\beta$ بحيث~$\OLId{greek-variable}{alpha} \in \beta \in \ordsucc{\OLId{greek-variable}{alpha}}$.
\end{enumerate}
\end{prop}

\begin{proof}
من التعريف مباشرة~$\OLId{greek-variable}{alpha} \in \OLId{greek-variable}{alpha} \cup \{\OLId{greek-variable}{alpha}\} = \ordsucc{\OLId{greek-variable}{alpha}}$.
ثم إن~$\ordsucc{\OLId{greek-variable}{alpha}}$ مجموعة متعدية من الأعداد الترتيبية، فهي
عدد ترتيبي بمقتضى~\olref[basic]{corordtransitiveord}. وأما
$\OLId{greek-variable}{alpha} \in \beta \in \ordsucc{\OLId{greek-variable}{alpha}}$ فمحال، إذ يلزم منه
$\beta \in \OLId{greek-variable}{alpha}$ أو~$\beta = \OLId{greek-variable}{alpha}$، وكلاهما يناقض
\olref[basic]{ordordered}.
\end{proof}
\noindent
ويتفرع عن ذلك وجه آخر للاستقراء المتناهي التجاوز.

\begin{thm}[الاستقراء المتناهي التجاوز البسيط]\ollabel{simpletransrecursion}
لتكن~$\phi(\OLId{latin-ordinary}{x})$ صيغة تتوافر فيها الشروط التالية:
\begin{enumerate}
\item $\phi(\emptyset)$؛
\item لكل عدد ترتيبي~$\OLId{greek-variable}{alpha}$، يلزم من~$\phi(\OLId{greek-variable}{alpha})$ أن~$\phi(\ordsucc{\OLId{greek-variable}{alpha}})$؛
\item إذا كان~$\OLId{greek-variable}{alpha}$ عددًا ترتيبيًا حديًا وكان
$(\forall \beta \in \OLId{greek-variable}{alpha})\phi(\beta)$، لزم~$\phi(\OLId{greek-variable}{alpha})$.
\end{enumerate}
إذن~$\forall \OLId{greek-variable}{alpha} \phi(\OLId{greek-variable}{alpha})$.
\end{thm}

\begin{proof}
نبرهن المقابل بالنقيض. فإن وجد عدد ترتيبي يحقق~$\lnot\phi$،
فليكن~$\gamma$ أصغره. إما~$\gamma = \emptyset$، وفيه مخالفة الشرط~\OLClassicalDigits{1}؛
وإما~$\gamma = \ordsucc{\OLId{greek-variable}{alpha}}$ لبعض~$\OLId{greek-variable}{alpha}$، فتقتضي أصغرية
$\gamma$ أن~$\phi(\OLId{greek-variable}{alpha})$، ويلزم من الشرط~\OLClassicalDigits{2} أن~$\phi(\gamma)$؛
وإما أن يكون~$\gamma$ حديًا، فتقتضي أصغريته
$(\forall \beta \in \gamma)\phi(\beta)$، ويلزم من الشرط~\OLClassicalDigits{3} أن~$\phi(\gamma)$.
ولا تخلو الحال من إحدى هذه الثلاث، وفي كل منها تناقض.
\end{proof}
\noindent
ونضيف ترميزًا أخيرًا نحتاج إليه فيما يأتي.

\begin{defn}\ollabel{defsupstrict}
لمجموعة الأعداد الترتيبية~$\OLId{latin-looped}{X}$ نضع
$\supstrict(\OLId{latin-looped}{X}) = \bigcup_{\OLId{greek-variable}{alpha} \in \OLId{latin-looped}{X}} \ordsucc{\OLId{greek-variable}{alpha}}$.
\end{defn}
\noindent
والرمز «lsub» اختصار لـ«أصغر حد علوي صارم».
\footnote{تكتب بعض الكتب «$\text{sup}(\OLId{latin-looped}{X})$» بهذا المعنى، وتكتب أخرى
«$\text{sup}(\OLId{latin-looped}{X})$» لأصغر حد علوي \emph{غير صارم}، أي~$\bigcup \OLId{latin-looped}{X}$.
ومتى كان لـ$\OLId{latin-looped}{X}$ عنصر أكبر~$\OLId{greek-variable}{alpha}$ اختلف المعنيان: فأصغر حد علوي
\emph{صارم} هو~$\ordsucc{\OLId{greek-variable}{alpha}}$، وأصغر حد علوي
\emph{غير صارم} هو~$\OLId{greek-variable}{alpha}$ نفسه.}
وبيان ذلك في القضية التالية.

\begin{prop}
إذا كانت~$\OLId{latin-looped}{X}$ مجموعة أعداد ترتيبية، فإن~$\supstrict(\OLId{latin-looped}{X})$ أصغر عدد
ترتيبي يتجاوز كل عدد ترتيبي في~$\OLId{latin-looped}{X}$.
\end{prop}

\begin{proof}
نضع~$\OLId{latin-looped}{Y} = \Setabs{\ordsucc{\OLId{greek-variable}{alpha}}}{\OLId{greek-variable}{alpha} \in \OLId{latin-looped}{X}}$، فيكون
$\supstrict(\OLId{latin-looped}{X}) = \bigcup \OLId{latin-looped}{Y}$. الأعداد الترتيبية متعدية، وكل عنصر
من أحدها عدد ترتيبي؛ ومن ثم فإن~$\supstrict(\OLId{latin-looped}{X})$ مجموعة متعدية من
الأعداد الترتيبية، فهي عدد ترتيبي بحسب~\olref[basic]{corordtransitiveord}.

فإذا كان~$\OLId{greek-variable}{alpha} \in \OLId{latin-looped}{X}$، كان~$\ordsucc{\OLId{greek-variable}{alpha}} \in \OLId{latin-looped}{Y}$، ولزم
$\ordsucc{\OLId{greek-variable}{alpha}} \subseteq \bigcup \OLId{latin-looped}{Y} = \supstrict(\OLId{latin-looped}{X})$، ثم
$\OLId{greek-variable}{alpha} \in \supstrict(\OLId{latin-looped}{X})$. إذن~$\supstrict(\OLId{latin-looped}{X})$ أكبر تمامًا من
كل عدد ترتيبي في~$\OLId{latin-looped}{X}$.

ولنأخذ~$\gamma$ حدًا علويًا صارمًا كيفما اتفق لـ$\OLId{latin-looped}{X}$؛ فيصح
$\OLId{greek-variable}{alpha} \in \gamma$ لكل~$\OLId{greek-variable}{alpha} \in \OLId{latin-looped}{X}$. وبما أن~$\gamma$ عدد
ترتيبي، يكون~$\ordsucc{\OLId{greek-variable}{alpha}} \subseteq \gamma$ لكل~$\OLId{greek-variable}{alpha} \in \OLId{latin-looped}{X}$.
ومن ثم~$\bigcup \OLId{latin-looped}{Y} = \supstrict(\OLId{latin-looped}{X}) \subseteq \gamma$، أي
$\supstrict(\OLId{latin-looped}{X}) \leq \gamma$. فثبت أن~$\supstrict(\OLId{latin-looped}{X})$ هو \emph{أصغر}
حد علوي صارم لـ$\OLId{latin-looped}{X}$.
\end{proof}

\end{document}
